% Authors: Nelson Lago, Arthur Del Esposte and Eduardo Zambom Santana
% Portions of the example contents: Arthur Del Esposte
% This file is distributed under the MIT Licence

%%%%%%%%%%%%%%%%%%%%%%%%%%%%%%%%%%%%%%%%%%%%%%%%%%%%%%%%%%%%%%%%%%%%%%%%%%%%%%%%
%%%%%%%%%%%%%%%%%%%%%%%%%%%%%%%%% PREÂMBULO %%%%%%%%%%%%%%%%%%%%%%%%%%%%%%%%%%%%
%%%%%%%%%%%%%%%%%%%%%%%%%%%%%%%%%%%%%%%%%%%%%%%%%%%%%%%%%%%%%%%%%%%%%%%%%%%%%%%%

% A língua padrão é a última citada
\documentclass[
  xcolor={hyperref,svgnames,x11names,table},
  hyperref={pdfencoding=unicode,plainpages=false,pdfpagelabels=true,breaklinks=true},
  brazilian,english,
]{beamer}

% Vários pacotes e opções de configuração genéricos
\usepackage{imegoodies}
\usepackage{caption}
\hidelinks % desabilita cor/sublinhado dos links (URLs, refs etc.)
\input{extras/source-code}


% Diretórios onde estão as figuras; com isso, não é preciso colocar o caminho
% completo em \includegraphics (e nem a extensão).
\graphicspath{{figuras/},{logos/}}

% Comandos rápidos para mudar de língua:
% \en -> muda para o inglês
% \br -> muda para o português
% \texten{blah} -> o texto "blah" é em inglês
% \textbr{blah} -> o texto "blah" é em português
\babeltags{br = brazilian, en = english}

% Espaçamento simples
\singlespacing


%%%%%%%%%%%%%%%%%%%%%%%%%%%% APARÊNCIA DO BEAMER %%%%%%%%%%%%%%%%%%%%%%%%%%%%%%%

%% Possible paper sizes: a0, a0b, a1, a2, a3, a4.
%% Possible orientations: portrait, landscape
%% Font sizes can be changed using the scale option.
\usepackage[size=a1,orientation=portrait,scale=1.8]{beamerposter}

\usetheme[poster]{imeusp} % carregado do diretório extras (veja basics.tex)

%Remove ícones de navegação
\beamertemplatenavigationsymbolsempty


%%%%%%%%%%%%%%%%%%%%%%%%%% COMANDOS PARA O USUÁRIO %%%%%%%%%%%%%%%%%%%%%%%%%%%%%

% Medidas feitas "a olho"
\newcommand\singlecol{\column{.963\textwidth}}
\newcommand\halfcol{\column{.46\textwidth}}
\newcommand\onethirdcol{\column{.2922\textwidth}}
\newcommand\twothirdscol{\column{.626\textwidth}}
\newcommand\onefourthcol{\column{.2084\textwidth}}

% Blocos de cor personalizada
\newenvironment{coloredblock}[2]%
  {
    \setbeamercolor{block title}{fg=white,bg=#1!80!white}
    \setbeamercolor{block body}{fg=darkgray,bg=#1!20!white}
    \setbeamercolor{local structure}{fg=darkgray,bg=#1!20!white}
    \begin{block}{#2}
  }
  {\end{block}}

% Bibliografia. Apenas estilos bibliográficos derivados de numeric,
% alphabetic, authortitle e authoryear (como beamer-ime) vão funcionar
% bem aqui! Outros estilos, como abnt ou apa, vão gerar problemas de
% layout que você vai precisar ajustar manualmente. Observe que, num
% poster ou apresentação, provavelmente é uma boa ideia usar apenas
% \nocite e não \cite.
\usepackage[
  %style=extras/plainnat-ime, % variante de autor-data, similar a plainnat
  %style=alphabetic, % similar a alpha
  %style=numeric, % comum em artigos
  %style=authoryear-comp, % autor-data "padrão" do biblatex
  %style=apa, % variante de autor-data, muito usado
  %style=abnt,
  style=extras/beamer-ime,
]{biblatex}

% Num poster, a bibliografia pode ficar em tamanho menor
\renewcommand*{\bibfont}{\footnotesize}


%%%%%%%%%%%%%%%%%%%%%%%%%%%%%%%%%%%%%%%%%%%%%%%%%%%%%%%%%%%%%%%%%%%%%%%%%%%%%%%%
%%%%%%%%%%%%%%%%%%%%%%%%%%%%%%%%%% METADADOS %%%%%%%%%%%%%%%%%%%%%%%%%%%%%%%%%%%
%%%%%%%%%%%%%%%%%%%%%%%%%%%%%%%%%%%%%%%%%%%%%%%%%%%%%%%%%%%%%%%%%%%%%%%%%%%%%%%%

% O arquivo com os dados bibliográficos para biblatex; você pode usar
% este comando mais de uma vez para acrescentar múltiplos arquivos
\addbibresource{bibliografia.bib}

% Este comando permite acrescentar itens à lista de referências sem incluir
% uma referência de fato no texto (pode ser usado em qualquer lugar do texto)
%\nocite{bronevetsky02,schmidt03:MSc, FSF:GNU-GPL, CORBA:spec, MenaChalco08}
% Com este comando, todos os itens do arquivo .bib são incluídos na lista
% de referências
%\nocite{*}

\title[CCSL]{An example poster using \LaTeX{}\\by CCSL-IME/USP}

\institute{Department of Computer Science --- University of São Paulo}

\author[ccsl@ime.usp.br]{\textbf{ccsl.ime.usp.br}}

\date{Month and day, year}

% Optional foot image
\footimage{\includegraphics[width=13cm]{ccsl-logo}}


%%%%%%%%%%%%%%%%%%%%%%%%%%%%%%%%%%%%%%%%%%%%%%%%%%%%%%%%%%%%%%%%%%%%%%%%%%%%%%%%
%%%%%%%%%%%%%%%%%%%%%%%%%%%%%%% INÍCIO DO POSTER %%%%%%%%%%%%%%%%%%%%%%%%%%%%%%%
%%%%%%%%%%%%%%%%%%%%%%%%%%%%%%%%%%%%%%%%%%%%%%%%%%%%%%%%%%%%%%%%%%%%%%%%%%%%%%%%

\begin{document}

% Em um poster não há \maketitle

\begin{frame}[fragile]\centering

\vspace{-.5\baselineskip}

%%%%%%%% Uma coluna %%%%%%%%
% (precisamos usar columns para alinhar com os demais blocos)
\begin{columns}[T]

  \singlecol

  \begin{block}{The CCSL logo (full width)}
    \centering
    \vspace{2\baselineskip} % Espaço adicional dentro do bloco (cor de fundo)
    \includegraphics[width=.9\textwidth]{ccsl-logo}
    \vspace{1\baselineskip}
  \end{block}

\end{columns}

    \vspace{1\baselineskip} % espaço adicional entre os blocos (branco)

%%%%%%%% Quatro colunas %%%%%%%%
\begin{columns}[T]

  \onefourthcol

  \begin{block}{One fourth width}
    \centering
    \vspace{.5\baselineskip}
    \includegraphics[width=.9\textwidth]{ccsl-logo}
    \vspace{.5\baselineskip}
  \end{block}

  \onefourthcol

  \begin{block}{One fourth width}
    \centering
    \vspace{.5\baselineskip}
    \includegraphics[width=.9\textwidth]{ccsl-logo}
    \vspace{.5\baselineskip}
  \end{block}

  \onefourthcol

  \begin{block}{One fourth width}
    \centering
    \vspace{.5\baselineskip}
    \includegraphics[width=.9\textwidth]{ccsl-logo}
    \vspace{.5\baselineskip}
  \end{block}

  \onefourthcol

  \begin{block}{One fourth width}
    \centering
    \vspace{.5\baselineskip}
    \includegraphics[width=.9\textwidth]{ccsl-logo}
    \vspace{.5\baselineskip}
  \end{block}

\end{columns}

\vspace{.5\baselineskip}

%%%%%%% Duas colunas %%%%%%%%
\begin{columns}[T]

  \halfcol

  \begin{block}{Pangrams (half width)}
    \begin{itemize}
      \item A \alert{pangram} is a sentence using every letter of a given
            alphabet at least once.
      \item Pangrams have been used to:
      \begin{itemize}
        \item display typefaces
        \item test equipment
        \item develop skills in handwriting, calligraphy, and keyboarding
      \end{itemize}
    \end{itemize}
  \end{block}

  \begin{block}{Examples of pangrams (half width)}
    \begin{itemize}
      \item In English
      \begin{itemize}
        \item The quick brown fox jumps over the lazy dog
        \item Sphinx of black quartz, judge my vow
        \item How vexingly quick daft zebras jump
        \item Pack my box with five dozen liquor jugs
      \end{itemize}
      \item In Portuguese
      \begin{itemize}
        \item Vejo xá gritando que fez show sem playback
        \item Já fiz vinho com toque de kiwi para belga sexy
        \item Dê já multa ao punk sexy que fez viação chegar à web
        \item Vejo galã sexy pôr quinze kiwis à força em baú achatado
      \end{itemize}
    \end{itemize}
  \end{block}

  \halfcol

  \begin{block}{The IME/USP logo (half width)}
    \centering
    % Ajustando o tamanho para alinhar o final com a coluna ao lado
    \vspace{.8\baselineskip}
    \includegraphics[width=.75\textwidth,trim=0 0 70 0,clip]{ime-logo}\par
    \vspace{.5\baselineskip}
  \end{block}

\end{columns}

\vspace{.5\baselineskip}

%%%%%%% Colunas assimétricas %%%%%%%%
  \begin{columns}[T]

  \twothirdscol

  \begin{block}{Bibliography (assimetrical columns -- two thirds width)}
    % "Estica" ligeiramente o bloco para alinhar com o final da outra coluna
    \vspace{.26\baselineskip}
    \nocite{FSF:GNU-GPL, MenaChalco08, biblatex}
    \printbibliography
    \vspace{.26\baselineskip}
  \end{block}

  \onethirdcol

  \begin{block}{A Table (one third width)}
  \begin{table}[H] % [H] é obrigatório com beamer!
    \centering
    \begin{tabular}{ccl}
      \toprule
      Code      & Abbreviation  & Name \\
      \midrule
      \texttt{A}  & Ala          & Alanine \\
      \texttt{C}  & Cys          & Cysteine \\
      \texttt{W}  & Trp          & Tryptophan \\
      \texttt{Y}  & Tyr          & Tyrosine \\
      \bottomrule
    \end{tabular}
  \end{table}

  \end{block}

\end{columns}

\vspace{.5\baselineskip}

%%%%%%%% Três colunas %%%%%%%%
\begin{columns}[T]

  \onethirdcol

    \begin{block}{Sponsors (one third width)}
    \vspace{.5\baselineskip}
    \centering
    \begin{figure}[H] % [H] é obrigatório com beamer!
      \begin{subfigure}[c]{.3\textwidth}
        \centering
        %\raisebox{.0972\baselineskip}{\includegraphics[height=.927\baselineskip]{fapesp-logo}}
        \includegraphics[width=.9\textwidth]{ccsl-logo}
      \end{subfigure}
      \begin{subfigure}[c]{.3\textwidth}
        \centering
        %\raisebox{-.53496\baselineskip}{\includegraphics[height=2.16\baselineskip]{capes-logo}}
        \includegraphics[width=.9\textwidth]{ccsl-logo}
      \end{subfigure}
      \begin{subfigure}[c]{.3\textwidth}
        \centering
        %\includegraphics[height=1.08\baselineskip]{cnpq-logo}
        \includegraphics[width=.9\textwidth]{ccsl-logo}
      \end{subfigure}
    \end{figure}
  \end{block}

  \onethirdcol

  \begin{block}{Sponsors (one third width)}
    \vspace{.5\baselineskip}
    \centering
    \begin{figure}[H] % [H] é obrigatório com beamer!
      \begin{subfigure}[c]{.3\textwidth}
        \centering
        %\raisebox{.0972\baselineskip}{\includegraphics[height=.927\baselineskip]{fapesp-logo}}
        \includegraphics[width=.9\textwidth]{ccsl-logo}
      \end{subfigure}
      \begin{subfigure}[c]{.3\textwidth}
        \centering
        %\raisebox{-.53496\baselineskip}{\includegraphics[height=2.16\baselineskip]{capes-logo}}
        \includegraphics[width=.9\textwidth]{ccsl-logo}
      \end{subfigure}
      \begin{subfigure}[c]{.3\textwidth}
        \centering
        %\includegraphics[height=1.08\baselineskip]{cnpq-logo}
        \includegraphics[width=.9\textwidth]{ccsl-logo}
      \end{subfigure}
    \end{figure}
  \end{block}

  \onethirdcol

  \begin{block}{Sponsors (one third width)}
    \vspace{.5\baselineskip}
    \centering
    \begin{figure}[H] % [H] é obrigatório com beamer!
      \begin{subfigure}[c]{.3\textwidth}
        \centering
        %\raisebox{.0972\baselineskip}{\includegraphics[height=.927\baselineskip]{fapesp-logo}}
        \includegraphics[width=.9\textwidth]{ccsl-logo}
      \end{subfigure}
      \begin{subfigure}[c]{.3\textwidth}
        \centering
        %\raisebox{-.53496\baselineskip}{\includegraphics[height=2.16\baselineskip]{capes-logo}}
        \includegraphics[width=.9\textwidth]{ccsl-logo}
      \end{subfigure}
      \begin{subfigure}[c]{.3\textwidth}
        \centering
        %\includegraphics[height=1.08\baselineskip]{cnpq-logo}
        \includegraphics[width=.9\textwidth]{ccsl-logo}
      \end{subfigure}
    \end{figure}
  \end{block}

\end{columns}

\end{frame}

\end{document}
